\begin{abstract}
Bài báo này đề xuất phương pháp mạng từ trở sinh lưới trong không gian ba chiều (3D-MBGRN) dựa trên hệ tọa độ trụ phục vụ phân tích trường điện từ phi tuyến của các cấu trúc máy điện quay đối xứng, với mô hình kiểm chứng là động cơ nam châm vĩnh cửu hướng trục (AFPM). Phương pháp thực hiện rời rạc hóa không gian tính toán dựa trên một cấu trúc lưới đồng nhất trong hệ tọa độ trụ $(r, \theta, z)$ để chia miền tính toán thành các phần tử hình lục diện đặc thù. Việc áp dụng linh hoạt các điều kiện biên chu kỳ và phản chu kỳ cho phép rút gọn hiệu quả miền tính toán dựa trên tính đối xứng hình học của máy điện. Cấu trúc lưới này cho phép mô tả chi tiết sự phân bố từ trường cục bộ tại mọi vị trí trong không gian máy điện, đồng thời kết hợp với kỹ thuật dịch chỉ số (index-shifting) để mô phỏng chuyển động quay của rotor một cách tự nhiên mà không cần tái tạo lưới hay nội suy ranh giới. Nhằm nâng cao tính ổn định số trị, bộ giải lặp điểm cố định được tích hợp thuật toán suy giảm thích nghi hệ số thư giãn vật liệu (Adaptive Material Decay) kết hợp cơ chế quay lui nghiệm (backtracking), giúp giảm thiểu rủi ro rơi vào bẫy trì trệ hội tụ. Kết quả kiểm chứng đối chiếu trực tiếp với phương pháp phần tử hữu hạn ba chiều (3D-FEM) ghi nhận sự tương đồng cao về dạng sóng transient và phổ hài của các đại lượng điện từ cốt lõi ở cả trạng thái không tải và có tải định mức, với độ lệch giữa hai phương pháp chỉ duy trì dưới $3.5\%$. Khảo sát hội tụ lưới cho thấy phương pháp 3D-MBGRN đạt sự hội tụ lưới với số lượng phần tử ít hơn 3D-FEM, đồng thời tiết kiệm đáng kể bộ nhớ RAM và cho tốc độ tính toán transient có thể nhanh hơn gấp đến $20$ lần tùy thuộc vào cấu hình lưới đối chiếu. Các kết quả đạt được thể hiện khả năng ứng dụng của phương pháp 3D-MBGRN như một công cụ tính toán hiệu quả thay thế cho 3D-FEM trong các quy trình thiết kế và tối ưu hóa máy điện.

\textbf{Keywords---} 3D Mesh-Based Reluctance Network (3D-MBGRN), Finite Element Method (FEM), Axial Flux Permanent Magnet Machine (AFPM), Periodic Boundary Conditions, Hexahedral Elements, Local Magnetic Field, Material Relaxation.
\end{abstract}